\section{関連研究}
\subsection{MissForestの概要}
missForest は，ランダムフォレストを基盤とした反復型の欠損値補完手法であり，分布仮定を必要としない
非パラメトリックな手法として提案されている．各変数を目的変数，残りの変数を説明変数とみなして
ランダムフォレストを学習し，欠損部分を予測する操作を変数ごとに繰り返すことで，データ行列全体の欠損値を
補完する枠組みを採用している．この反復的な学習・更新により，特徴量間の非線形な関係や相互作用を考慮した
補完が可能である．

missForest の特徴として，連続変数とカテゴリ変数が混在するデータに適用可能である点や，補完精度の推定に
ランダムフォレストの OOB誤差を利用できる点が挙げられる．OOB 誤差を用いることで，
外部の検証用データを用いることなく，補完結果の誤差を近似的に評価できる．また，多数の決定木の予測結果を
平均化するランダムフォレストの構造は，多重代入的な性質を内包していると整理されている．

原著論文およびその後の研究において，missForest は医療データや生物データをはじめとする実データを用いた
実験により，他の欠損値補完手法と比較して安定した補完性能を示すことが報告されており，機械学習に基づく
欠損値補完手法の代表的な手法として位置づけられている．一方で，高次元データにおいては，補完に寄与しない
特徴量の影響を受けやすいことが指摘されており，特徴量の扱いに関して改良の余地があるとされている．
